%% ============================================================
%% 片段：s04b-lines   对应原文 PDF 第 98--109 页
%% 本文件由 gw/convert.py 从 arXiv 源码（hep-th/0612073v2）自动生成：
%%   * 公式/编号/引用/图 已经锁定为原文内容，请勿改动任何数学内容与编号；
%%   * 只需把英文正文译为中文，并在每节开头加“本节重点”、关键处加译注。
%% ============================================================


\subsection{线算符与单值性}
%% ==== 原文 §4.5：PDF 第 98 页（印刷第 97 页）====
\label{sec:4.5}

在没有分歧时，支撑在点
$q\in C$ 处（再乘以 $\Sigma$ 中沿其边界延伸的一条线）的线算符（line operator）
自然地作用在膜上。这一点在
文献~\cite{ref8} 第 6.4 节有解释。

在 $\Psi=\infty$ 处，自然的线算符是威尔逊算符（Wilson operator），
它们由规范群的一个表示的选择来分类。我们把 $\Psi=\infty$ 处的规范群
取为 $^L\neg
G$，于是威尔逊算符由该群的表示来标记。一般地，威尔逊算符
可以改变离散电通量 ${\bf e}_0$。保持 ${\bf e}_0$ 不变的
威尔逊算符，是那些来自 $^L\neg G_\ad$（即 $^L\neg G$ 的伴随形式）的表示的算符。

在 $\Psi=0$ 处，自然的线算符是 't Hooft 算符，
规范群为 $G$。$G$ 规范场论中的 't Hooft 算符是用
奇异的 BPS 磁单极子构造出来的，同样
由 $^L\neg G$ 的表示分类。一般地，'t Hooft
算符改变离散磁通量 ${\bf m}_0$（在对偶下 ${\bf e}_0$ 变换为它）。
保持 ${\bf m}_0$ 不变的 't Hooft
算符，是那些与 $^L\neg G_{\rm ad}$ 的
表示相联系的算符。

所以在 $\Psi=\infty$ 或 $\Psi=0$ 处，在没有分歧时
作用于给定点 $q\in C$ 的线算符，
对应于 $^L\neg G$ 的表示；若只想要保持离散通量不变的算符，则
对应于 $^L\neg G_\ad$ 的表示。这些威尔逊算符或 't Hooft 算符的
复合对应于表示的张量积。几何朗兰兹纲领的一个核心论断是：
$^L\neg G$ 与 $G$ 之间的对偶把一个由 $^L\neg G$ 的表示标记的威尔逊算符
映为一个由同一表示标记的 't Hooft 算符。

现在设 $p\in C$ 是一个分歧点。我们要确定
$p$ 处究竟是什么结构取代了一般点处威尔逊
或 't Hooft 线算符的作用。我们将证明，这个问题的答案可以用
分歧参数空间中的单值性来描述。

\bigskip\noindent{\it 对偶对称的单值性}

我们先以对偶对称的方式重复第~\ref{sec:3.7}~节的推理。为简单起见，
从列维子群为 $\mathbb{L}=\TT$ 的情形开始。

对规范群 $G$，$p$ 点处带分歧结构的定义依赖于参数
$(\alpha,\beta,\gamma,\eta)\in \mathfrak t^3\times \mathfrak t^\vee$
的选取，并模去某个等价群的作用。这个群
由如下变换生成：{\it (i)} 用格 $\Lambda_\cochar$ 平移 $\alpha$；
{\it (ii)} 用对偶格 $\Lambda_\char$ 平移 $\eta$；
以及 {\it (iii)} 外尔群 $\Weyl$。所以作用在参数上的群是一个
扩张
% 式 (4.11)
\setcounter{equation}{10}
\begin{equation}\label{eq:4.11}
  \hat {\mathcal
V}=\left(\Lambda_\cochar\oplus\Lambda_\char\right)\rtimes\Weyl.
\end{equation}

商 $(\mathfrak t^3\times \mathfrak t^\vee)/\hat{\mathcal V}$ 与
$(\TT\times \mathfrak t\times\mathfrak t\times ^L\neg\TT)/\Weyl$ 相同。

若只想保留不改变离散通量 ${\bf m}_0$ 与 ${\bf e}_0$ 的
$\alpha$、$\eta$ 平移，则应把 $\Lambda_\cochar$ 换成 $\Lambda_\cort$、
把 $\Lambda_\char$ 换成 $\Lambda_\rt$。于是得到一个更小的群

% 式 (4.12)
\setcounter{equation}{11}
\begin{equation}\label{eq:4.12}
  {\mathcal
V}=\left(\Lambda_\cort\oplus\Lambda_\rt\right)\rtimes\Weyl.
\end{equation}

${\mathcal V}$ 可说是仿射外尔群 $\AffWeyl$ 的
一种对偶对称版本，因为 $G$ 的仿射外尔群是
$\Lambda_\cort\rtimes \Weyl$，而 $^LG$ 的仿射外尔群
是 $\Lambda_\rt\rtimes \Weyl$。${\mathcal
V}$ 与 $\hat {\mathcal V}$ 之间的关系极其简单。商
$\Lambda_\cochar/\Lambda_\cort$ 就是 $\pi_1(G)$，同样
$\Lambda_\char/\Lambda_\rt=\pi({}^L\neg G)$。外尔群在这些群上
平凡地作用，所以我们有一个群扩张

% 式 (4.13)
\setcounter{equation}{12}
\begin{equation}\label{eq:4.13}
  1\to {\mathcal V}\to \hat{\mathcal V}\to \pi_1( G)\times
\pi_1({}^L\neg G)\to 1.
\end{equation}

\begin{yizhu}[原文疑误]上式前一式中 $\Lambda_\cochar/\Lambda_\cort=\pi_1(G)$，与此对应，
这里的 $\pi({}^L\neg G)$ 显然应为 $\pi_1({}^L\neg G)$（即 $^L\neg G$ 的基本群；群同态也正需要
$\pi_1({}^L\neg G)$ 才与 (4.13) 中的 $\pi_1({}^L\neg G)$ 相符）。此处照原文保留，建议核对 PDF 第 99 页。\end{yizhu}

\begin{yizhu}[易错点]$\hat{\mathcal V}$ 与 $\mathcal V$ 的差别正是离散通量：由 (4.13)，
$\hat{\mathcal V}/\mathcal V\cong\pi_1(G)\times\pi_1({}^L\neg G)$，其中 $\pi_1(G)$ 由 $\Lambda_\cochar/\Lambda_\cort$
实现（对应离散磁通量 ${\bf m}_0$），$\pi_1({}^L\neg G)$ 由 $\Lambda_\char/\Lambda_\rt$ 实现（对应离散电通量
${\bf e}_0$）。因此“保持通量不变”的对称性要用 $\mathcal V$ 描述，而含改变通量的完整对偶对称性要用
$\hat{\mathcal V}$；由于外尔群在商上平凡作用，此扩张实为阿贝尔扩张。\end{yizhu}

我们说点 $(\alpha,\beta,\gamma,\eta)\in \mathfrak
t^3\times \mathfrak t^\vee$ 是 ${\mathcal V}$ 正则的，如果它不被
$\mathcal V$ 的任何非平凡元素固定。根据第~\ref{sec:3.6}~节，
按此意义不正则的点，正是以
$\MH(\alpha,\beta,\gamma,\eta;p)$ 为靶空间的西格玛模型出现局部奇点的点。我们令
$\eury$ 表示 $\mathfrak t^3\times\mathfrak t^\vee$ 中非 ${\mathcal V}$ 正则点集的补集。
（若有不止一个分歧点，定义 $\eury$ 时还要去掉整体奇点的轨迹；
它具实余维四，不影响论证。）群 ${\mathcal V}$ 自由地作用在 $\eury$ 上，所以
商 $\eury/{\mathcal V}$ 是光滑流形。

于是我们得到一族以流形 $\eury/{\mathcal V}$ 为参数的光滑 $(4,4)$ 西格玛模型
（还依赖其他数据，例如规范耦合参数 $\tau$）。现在可以作一个论证，
它可说是第~\ref{sec:3.7}~节推理的量子版本。关于膜的一个基本事实是：
每个膜都有一个荷，即一个 $K$-理论类，粗略地说
取值于靶空间 $\MH$ 的复 $K$-理论中。$K(\MH)$ 随
$(\alpha,\beta,\gamma,\eta)$ 变化，表现为
$\eury/{\mathcal V}$ 上某个平坦丛的纤维。取单值性，便得到 $\mathcal V$
在 $K(\MH)$ 上的一个作用。

要把这一点说得更精确，必须回忆：$\MH$ 有
由 ${\bf m}_0$ 标记的不同拓扑分支。此外，对每个 ${\bf e}_0$，
可定义一个扭曲膜的空间，其荷取值于
$K$-理论的一个扭曲版本\footnote{这就是扭曲向量丛的
$K$-理论，它在当前语境下的描述见文献~\cite{ref8} 第 7.1 节。其想法是：
${\bf e}_0$ 定义了 $\MH$ 上的一个平坦 gerbe，而 $K_{{\bf e}_0}(\MH)$ 是相对于该 gerbe
定义的扭曲 $K$-理论。}
 中，我们不妨称之为 $K_{{\bf e}_0}(\MH)$。设 $K_{{\bf e}_0}(\MH;{\bf m}_0)$ 为
$\MH$ 中由 ${\bf m}_0$ 定义的分支的 ${\bf e}_0$ 扭曲 $K$-理论。上述论证表明，
对每个 ${\bf e}_0$ 与 ${\bf m}_0$ 的选取，${\mathcal V}$
都作用在 $K_{{\bf e}_0}(\MH;{\bf m}_0)$ 上。更大的群 $\hat {\mathcal V}$ 也
作用，只须把上述推理稍作推广；当然，它的作用
置换 ${\bf e}_0$ 与 ${\bf
m}_0$ 的可能取值。特别地，作为阿贝尔群，$K_{{\bf
e}_0}(\MH;{\bf m}_0)$ 在同构意义下与 ${\bf e}_0$、${\bf
m}_0$ 无关。

我们可以用比较具体的办法理解 ${\mathcal V}$（或其扩张
$\hat{\mathcal V}$）如何作用在 $K(\MH)$ 上。子群
${\mathcal W}_{\rm aff}=\Lambda_{\cort}\rtimes\Weyl$ 的作用
已在第~\ref{sec:3.7}~节描述过了。

另一方面，平移 $\eta\to \eta+v$（$v\in
\Lambda_\rt$）是世界面 $B$ 场的平移，它使
其周期改变 $2\pi$ 的整数倍。
这样的平移作用在膜 $\mathcal
B$ 上，就是把 $\mathcal B$ 的 Chan-Paton 向量丛
与第一陈类为 $v$ 的线丛 ${\mathcal L}^v$ 作张量积。于是对 $K$-理论
的相应作用，就只是与 ${\mathcal L}^v$ 作张量积，$v\in \Lambda_\rt$。
由于这一作用在经典层面是清楚的，
$\mathcal V$ 作用在 $K(\MH)$ 上这一事实本身
并不能告诉我们多少超出 $\AffWeyl$ 作用之外的信息，
而后者我们已经从第~\ref{sec:3.7}~节知道了。

然而，看起来能给出更多信息的，是
$S:\tau\to -1/n_{\mathfrak g}\tau$ 的作用。它把 $K(\MH(^L\neg G,C))$ 与
$K(\MH(G,C))$ 互换，同时把 $\Lambda_{\rm rt}$ 与
$\Lambda_\cort$（以及 ${\bf e}_0$ 与 ${\bf m}_0$）互换。因此它
把“$\Lambda_{\rm rt}$ 通过线丛扭转的作用”与“$\Lambda_\cort$ 通过单值性的作用”
互换。$\MH$ 的 $K$-理论的这一对称性在经典层面并不显然。

我们简要说明一个说明其运作的简单例子。$K(\MH)$ 的乘法结构
的单位元，是某个膜 ${\mathcal B}$ 的 $K$-理论类：该膜的支撑就是整个
$\MH$，其 Chan-Paton 线丛平凡。这个膜在
我们改变 $\alpha$、$\beta$、$\gamma$ 时连续变化，
所以它的 $K$-理论类在
仿射外尔群 $\AffWeyl=\Lambda_\cort\rtimes\Weyl$ 的作用下不变。于是
对偶操作 $S$ 必须把它变成一个在 $\Lambda_\rt$ 作用下
不变的膜 $\tilde{\mathcal
B}$。换句话说，$\tilde {\mathcal B}$ 的 $K$-理论类必须在
如下操作下不变：把 $\tilde
{\mathcal B}$ 的 Chan-Paton 丛与形如 ${\mathcal L}^v$ 的线丛作张量积。这一
论断意味着：限制到 $\tilde{\mathcal B}$ 的支撑上时，${\mathcal L}^v$ 必须（拓扑上）平凡。
实际上，$\tilde{\mathcal B}$ 是一个支撑在 Hitchin
纤维化某个截面上的膜，所以它的支撑
同构于该纤维化的底空间 $\EUBB$。由于
这是一个可缩空间（同构于 $\C^N$，$N$ 为某个整数），
任何线丛 ${\mathcal L}^v$ 限制到
$\tilde{\mathcal B}$ 的支撑上时确实平凡。

\bigskip\noindent{\it $A$-模型的单值性}

迄今为止所讲的内容应当说是有趣的，但看来并不能使我们
很接近威尔逊与 't Hooft 算符的类比物。

要更进一步，必须讨论的不是一般的膜，而是特定种类的膜。
例如，我们来讨论规范群为 $G$ 的 $K$ 型 $A$-模型的膜。这是
$\Psi=0$ 处相关的模型。该模型的定义依赖于
参数 $\gamma$ 与 $\eta$。所以要研究
$A$-模型的膜，应当保持 $\gamma$ 与 $\eta$ 固定。另一方面，
$A$-模型局部地与 $\alpha$、$\beta$ 无关。
因此，通过改变 $\alpha$ 与 $\beta$ 得到的单值性群
将是 $A$-模型的一个对称群。这个群
依赖于 $\gamma$ 与 $\eta$。

在保持 $\MH$ 作为辛簇（辛结构为 $\omega_K$）不变的同时改变
$\alpha$ 与 $\beta$ 的过程，已在第~\ref{sec:3.2}~节讨论过。
由该讨论可知，当在 $\alpha-\beta$ 空间中绕一个回路一周
（避开奇点并保持 $\gamma$ 固定）时，$\MH$ 会通过一个辛同胚（symplectomorphism）改变。
因此，我们即将分析的 $A$-模型的对称群
就是一个经典辛同胚群（准确地说，是辛同胚群的
一个分支群）。这与 $B$-模型的对偶单值性形成对照：
后者涉及改变量子参数 $\eta$，
因而只在量子意义下才有意义。

现在讨论 $\mathcal V$ 的作用。由于 $\Lambda_\rt$
只作用于 $\eta$，而在研究 $A$-模型及其单值性时 $\eta$ 是固定的，
所以它不会起什么有趣的作用。我们不妨
先模去 $\Lambda_\rt$，并把 $\gamma,\eta$ 这一对换成
$Y=\exp(-2\pi(\eta-i\gamma))$。粗略地说，$Y$ 取值于
$^L\neg G_\C$ 的复极大环面 $^L\TT_\C$。
更精确地说，由于我们只模去了 $\Lambda_\rt$ 而非
$\Lambda_\char$，$Y$ 取值于 $^L\bar
G_\C$ 的极大环面，而 $^L\bar G_\C$ 是 $^L\neg G_\C$ 的单连通覆盖群。
由于 $A$-模型依赖于 $Y$，我们不关心涉及改变 $Y$ 的
单值性。所以外尔群中相关的部分是保持 $Y$ 不变的
子群。我们把它记作 ${\mathcal W}_Y$。


$\mathcal V$ 中只作用于 $\alpha$ 与 $\beta$ 的子群，
是 $\Weyl_Y$ 被 $\Lambda_\cort$ 所作的扩张：

% 式 (4.14)
\setcounter{equation}{13}
\begin{equation}\label{eq:4.14}
  {\mathcal V}_Y=\Lambda_\cort\rtimes \Weyl_Y.
\end{equation}

 于是若 $Y=1$，${\mathcal V}_Y$ 就是仿射外尔群。我们
说序对 $(\alpha,\beta)\in \mathfrak t\times \mathfrak t$ 是
${\mathcal V}_Y$ 正则的，如果它不被 ${\mathcal V}_Y$ 中除恒等元以外的
任何元素固定。非 ${\mathcal V}_Y$ 正则的序对
对应于不正则的四元组 $(\alpha,\beta,\gamma,\eta)$，
从而对应于 $\MH$ 的局部奇点。我们定义
$\eurz_Y$ 为 ${\mathcal V}_Y$ 正则序对构成的空间。群 ${\mathcal
V}_Y$ 自由地作用在 $\eurz_Y$ 上，于是得到一族以 $\eurz_Y/{\mathcal V}_Y$ 为参数的光滑西格玛模型。
若从西格玛模型过渡到相应的 $K$ 型 $A$-模型，这族模型
变成局部常值的，因为这个 $A$-模型局部地与参数 $\alpha$ 和 $\beta$ 无关。
所以基本群 $B_Y=\pi_1(\eurz_Y/{\mathcal V}_Y)$ 作为
$A$-模型的自同构群作用。（$B$ 大致代表“辫”（braid），
稍后会更清楚。）

这一推理与我们先前讨论
$\AffWeyl$ 与 $\mathcal V$ 的作用时所讲的类似。
不过，有两个值得注意的差别：

{\it (A)} 非 ${\mathcal V}_Y$ 正则的序对 $(\alpha,\beta)$
一般具实余维二，原因很简单：我们只考虑
两个变量 $\alpha$ 与 $\beta$。（相比之下，
不正则的三元组 $(\alpha,\beta,\gamma)$ 或四元组
$(\alpha,\beta,\gamma,\eta)$ 分别具实余维三或
四。）因此，${\mathcal V}_Y$ 正则序对构成的空间 $\eurz_Y$
未必单连通，商 $\eurz_Y/{\mathcal V}_Y$ 的基本群
$B_Y$ 一般也不等于 ${\mathcal V}_Y$。不过我们确实得到
一个正合列

% 式 (4.15)
\setcounter{equation}{14}
\begin{equation}\label{eq:4.15}
  1\to \pi_1(\eurz_Y)\to B_Y\to {\mathcal V}_Y\to
1.
\end{equation}


{\it (B)} 与此相关，在两或更多分歧点的情形，
整体奇点具实余维二，在完整的处理中
可能起作用。若只考虑支撑于这样的希格斯丛
$(E,\varphi)$ 上的膜——其底丛 $E$ 在忘掉
抛物结构后是稳定的——便可避开它们。在本文中，
我们不关心整体奇点。


另一种有用的描述是先模去 $\Lambda_\cort$ 对 $\alpha$ 的作用。
模去这一作用后，
序对 $(\alpha,\beta)$ 合并为元素 $T=\exp(-2\pi(\alpha+i\beta))\in \bar\TT_\C$，
即 $G_\C$ 的单连通覆盖 $\bar G_\C$ 的极大
环面。我们说 $T$ 是 ${\mathcal W}_Y$ 正则的，如果它不被
${\mathcal W}_Y$ 的任何非平凡元素固定。记 $\eurq_Y$ 为
$\bar\TT_\C$ 中 ${\mathcal W}_Y$ 正则元素构成的空间。于是 ${\mathcal W}_Y$
自由地作用在 $\eurq_Y$ 上，且 $B_Y=\pi_1(\eurq_Y/\Weyl_Y)$。

照例，为了对群 $B_Y$ 有些理解，
考察 $Y$ 正则或 $Y=1$ 的情形是有益的。其他
情形介于这两者之间。在极端情形下，我们有：

(1) 若 $Y$ 正则，则 $\Weyl_Y=1$，且 $\eurq_Y=\bar\TT_\C$。
于是 $B_Y=\pi_1(\eurq_Y/\Weyl_Y)=\pi_1(\eurq_Y)=\Lambda_\cort$。

(2) 另一种极端是 $Y=1$，此时 $\Weyl_Y=\Weyl$。于是 $\eurq_Y$
是 $\bar\TT_\C$ 中正则点构成的空间，
而 $\eurq_Y/\Weyl_Y=\eurq_Y/\Weyl$ 是 $\bar G_\C$ 中正则共轭类构成的空间。
这个空间的基本群称为 $^L\neg G$ 的仿射辫群\footnote{“仿射辫群”这一名称的
动因如下。$n$ 个字母上的普通辫群，是 $\C$ 中 $n$ 个互异点的无序
$n$ 元组构成的空间的基本群。若把 $\C$ 换成
$\C^*$，便得到 $GL(n,\C)$ 的仿射辫群，因为 $\C^*$ 中的一个无序
$n$ 元组等价于 $GL(n,\C)$ 的一个正则半单共轭类。}。
我们把它记作
$B_{\rm aff}(^L\neg G)$。于是 $Y=1$ 时的单值性群为 $B_1=
B_{\rm aff}(^L\neg G)$。现在，当 $Y=1$ 时，${\mathcal
V}_Y=\Lambda_\cort\rtimes \Weyl$ 就是 $^L\neg G$ 的仿射外尔群。
此时式~\eqref{eq:4.15} 给出正合列

% 式 (4.16)
\setcounter{equation}{15}
\begin{equation}\label{eq:4.16}
  1\to \pi_1(\eurz)\to B_{\rm aff}(^L\neg G)\to
\AffWeyl(^L\neg G)\to 1,
\end{equation}
 其中 $\eurz=\eurz_1$ 是
正则序对 $(\alpha,\beta)\in\bar\TT\times \mathfrak t$ 构成的空间，或等价地，
是 $\AffWeyl$ 正则序对 $(\alpha,\beta)\in\mathfrak
t\times \mathfrak t$ 构成的空间。

仿射辫群是在文献~\cite{ref29} 中引入的，其背景大致上是
本问题的一个局部版本。我们已在第~\ref{sec:3.8}~节讨论过二者的关系，
并将在第~\ref{sec:6}~节回到这一点。

我们定义的群 $B_Y$，是 $A$-模型的分歧参数中保持
${\bf m}_0$ 不变的那些单值变换构成的群。我们可以放宽这一条件，
考虑一个扩张群 $\hat B_Y=B_Y\rtimes {\mathcal Z}(G)$，
它包含改变 ${\bf m}_0$ 的变换。对正则的 $Y$，
$\hat{B}_Y=\Lambda_\cochar$；而对 $Y=1$，它是一个扩张

% 式 (4.17)
\setcounter{equation}{16}
\begin{equation}\label{eq:4.17}
  \hat{B}_{\rm aff}({}^L\neg G)= B_{\rm aff}(^L\neg
G)\rtimes {\mathcal Z}(G).
\end{equation}


以上分析是在 $\Psi=0$ 处进行的。显然，对复结构为 $J$
的 $B$-模型可以作类似分析，后者对应于 $\Psi=\infty$。这会把
$\alpha$ 与 $\eta$ 的角色以及 $G$ 与 $^LG$ 的角色互换。净效果是：
作用在规范群为 $G$ 的 $K$ 型 $A$-模型的膜上的单值性群，
与规范群为 $^L\neg G$ 的 $J$ 型 $B$-模型的单值性群相同。
（不过，正如上面指出的，$A$-模型的单值性群通过经典辛同胚作用，
而 $B$-模型的单值性群则高度非经典。）正如我们下面将要讨论的，
这推广了一种情形下的 't Hooft 算符与另一种情形下的威尔逊算符之间的对应。

\bigskip\noindent{\it 与线算符的关系}

我们把这一切都用作用在 $\Psi=\infty$ 的 $B$-膜或 $\Psi=0$ 的
$A$-膜上的群 ${\mathcal V}_Y$ 来表述。
更一般地，由于可以取两个膜的直和，或
用一个正整数乘一个膜，$\mathcal V_Y$ 的元素的
正整数线性组合可以作用在膜上。这样的线性组合构成
我们所谓的 $\mathcal V_Y$ 的群半环（group semiring）。

这里我们主要考虑 $Y$ 正则的情形。此时 ${\mathcal
V}_Y=\Lambda_\cort$，而 $\hat {\mathcal V}_Y=\Lambda_\cochar$。设
${\mathcal B}$ 是一个 Chan-Paton 丛为 ${\mathcal U}$ 的膜。我们在
上面发现，$\Lambda_\cort$ 与 $\Lambda_\cochar$ 通过 ${\mathcal U}\to {\mathcal L}^v\otimes {\mathcal U}$
作用在 ${\mathcal B}$ 上，其中 $v\in\Lambda_\cort$ 或 $v\in\Lambda_\cochar$。所以群半环
通过与此类线丛的直和作张量积来作用：

% 式 (4.18)
\setcounter{equation}{17}
\begin{equation}\label{eq:4.18}
  {\mathcal U}\to \left(\oplus_{i=1}^w{\mathcal
L}^{v_i}\right)\otimes {\mathcal U}.
\end{equation}


下面我们要证明，在 $Y$ 正则的情形，这一结构究竟
如何与无分歧时存在的线算符相联系。我们
用 $B$-模型的语言来表述。（背景见文献~\cite{ref8} 第 7、8 节。）设
$^L\neg R$ 是 $^L\neg G$ 的一个表示。设 $({\mathcal
E},\hat\varphi)$ 是 $\MH(^L\neg
G,C)\times C$ 上的万有希格斯丛（universal Higgs bundle），并设 ${\mathcal E}_{^L\neg R}$ 是该丛在表示
$^L\neg R$ 中的相伴丛\footnote{若 $^L\neg G$ 的中心在表示 $^L\neg R$ 中非平凡地作用，
则 ${\mathcal E}_{^L\neg R}$ 是扭曲向量丛而非普通向量丛，
相应的威尔逊算符会改变 ${\bf e}_0$。}。记 ${\mathcal E}_{^L\neg R}|_q$ 为
${\mathcal E}_{^L\neg R}$ 在 $\MH\times q$（$q$ 为 $C$ 中的点）上的限制。
现在考虑 $q$ 处表示 $^L\neg R$ 中的威尔逊算符 $W_q(^L\neg R)$。
它对膜 ${\mathcal B}$ 的作用如下：若 ${\mathcal U}$ 是膜 ${\mathcal B}$
的 Chan-Paton 向量丛，则用 $W_q(^L\neg R)$ 作用于 ${\mathcal B}$ 所得的
膜的 Chan-Paton 丛为 ${\mathcal E}_{^L\neg R}|_q\otimes {\mathcal
U}$。

现在让 $q$ 趋近一个使 $Y$ 正则的分歧点 $p$。
在 $p$ 处，$^L\neg G$-丛 $E$ 的结构群
约化为极大环面 $^L\TT$，相应地，
$^L\neg G$ 丛 $E\to C$ 在 $p$ 处的纤维分解为
与 $^L\TT$ 的表示相对应的子空间的直和。
于是对 $q=p$，向量丛 ${\mathcal E}_{^L\neg R}|_q$ 分解为
线丛 ${\mathcal L}^v$（$v\in \Lambda_\char$）的直和。与这种线丛之直和
作张量积，正是群半环中某个元素作用的一个例子，
如式~\eqref{eq:4.18} 所述。

这样就证明了：至少在 $Y$ 正则时，我们从一般点 $q$ 处的
威尔逊或 't Hooft 算符所得到的，在 $q$ 趋近歧点 $p$ 的极限下，
是 $B$-模型或 $A$-模型的单值性所给出结果的特殊情况。
正合列~\eqref{eq:4.15} 表明，要把这一点推广到非正则的 $Y$ 是棘手的。
威尔逊线算符 $W(^L\neg R)$ 与我们上面所见类似，
给出 $\mathcal V_Y=\Lambda_\char({}^L\neg G)\rtimes
{\mathcal W}_Y$ 的群半环中的一个元素。但把它提升为 $B_Y$ 的群半环中
的元素并无唯一方式。事实上，$q$ 趋近 $p$ 的“方向”
（在黎曼面 $C$ 上）会影响威尔逊算符的极限，
这一点已在另一种语言中分析过~\cite{ref75}。

我们给出的解释可能显得相当抽象。在第~\ref{sec:5}~节，我们将
在规范场论中更直接地重新考察这些问题。对 $A$-模型我们将
得到相当令人满意的描述，对 $B$-模型则能得到一些洞见。

\bigskip\noindent{\it 对任意 $\mathbb{L}$ 的推广}

以上一切都可以自然推广到 $G$ 的任意列维子群 $\mathbb{L}$。
要做的只是把 $\mathbb{L}$ 纳入所有陈述之中。

当然，首先我们要求参数
$(\alpha,\beta,\gamma,\eta)$ 在外尔群
${\mathcal W}_\mathbb{L}$ 下不变。为简单起见集中讨论
$A$-模型：为描述其单值性，必须保持
$Y=\exp(-2\pi(\eta-i\gamma))$ 固定。我们记
${\mathcal W}_{Y,\mathbb{L}}$ 为 $G$ 的外尔群中
使 ${\mathcal W}_\mathbb{L}$ 正规化并与 $Y$ 交换的子群。这正是
作用在 $(\alpha,\beta)$ 上、保持其
$\mathbb{L}$ 不变性并保持 $Y$ 不变的子群。在平移
$\alpha$ 时，我们应限于 ${\mathcal
W}_\mathbb{L}$ 不变的子格 $\Lambda_\cort^\mathbb{L}\subset
\Lambda_\cort$。所以式~\eqref{eq:4.14} 的类比是：只作用于 $\alpha$ 与 $\beta$ 的群
现在是

% 式 (4.19)
\setcounter{equation}{18}
\begin{equation}\label{eq:4.19}
  {\mathcal
V}_{Y,\mathbb{L}}=\Lambda_\cort^\mathbb{L}\rtimes{\mathcal W}_{Y,\mathbb{L}}.
\end{equation}


记 $\eurz_{Y,\mathbb{L}}$ 为 $\mathbb{L}$ 不变的序对
$(\alpha,\beta)\in \mathfrak t\times \mathfrak
t$ 中不被 ${\mathcal V}_{Y,\mathbb{L}}$ 的任何元素固定的那些构成的空间。群 ${\mathcal V}_{Y,\mathbb{L}}$
自由地作用在 $\eurz_{Y,\mathbb{L}}$ 上，于是得到一族以商 $\eurz_{Y,\mathbb{L}}/{\mathcal V}_{Y,\mathbb{L}}$ 为参数的
光滑 $A$-模型。这族模型是局部常值的，
所以基本群 $B_{Y,\mathbb{L}}=\pi_1(\eurz_{Y,\mathbb{L}}/{\mathcal V}_{Y,\mathbb{L}})$
作为
$A$-模型的自同构群作用。这一陈述对 $B$-模型有直接的类比。

\subsection{表示与膜}

%% ==== 原文 §4.6：PDF 第 106 页（印刷第 105 页）====
\label{sec:4.6}

在没有分歧时，一个不可约的平坦 $^LG$-丛 $E\to C$ 对应于
$\MH$ 中的一个光滑点 $x_E$。$\MH$ 的 $B$-模型与 Kähler 度量无关，
而后者由规范耦合参数 $\tau=\theta/2\pi+4\pi i/e^2$ 控制。
取 ${\rm Im}\,\tau$ 充分大，便可进入 $B$-模型可以
半经典处理的区域。此时存在一个支撑在
光滑点 $x_E\in\MH$ 上的零膜（zerobrane）${\mathcal B}_E$。它是威尔逊线算符作用的
本征膜（eigenbrane），也是几何朗兰兹纲领中的
首要研究对象。

这里 $E$ 的不可约性所起的作用，是使我们远离
$\MH$ 的奇点。若 $E$ 可约（但
半稳定），它仍对应于 $\MH$ 中的一点 $x_E$，只是该点
奇异。我们仍可考虑支撑在 $x_E$ 处的膜，
但这类膜的理论更复杂。一般说来，
并不存在与 $E$ 相联系的唯一且典范确定的
膜。取而代之，我们可以定义一个支撑在 $x_E$ 处的膜的
空间或“范畴”；可以论证这个空间在威尔逊算符的作用下
封闭。这一复杂性是数论中一个已知现象的
几何类比。

我们想讨论带分歧情形的类似问题，即
$C\backslash p$ 上的平坦丛
（更一般地，对某个有限集 $\{p_1,\dots,p_s\}$，$C\backslash\{p_1,\dots,p_s\}$ 上的平坦丛），
它绕 $p$ 有单值性。设 $E$ 是这样的丛，$V$ 表示它绕 $p$ 的单值性。我们假定
$E$ 不可约，否则就会像无分歧时那样遇到整体奇点。
这里的目标是考察只依赖于 $V$ 的共轭类的局部奇点
所起的作用。

对任意 $V$，我们取 $^L\neg\alpha$ 与 $^L\neg\gamma$，使得 $U=\exp(-2\pi({}^L\neg\alpha-i
{}^L\neg\gamma))$
落在 $V$ 的轨道的闭包中。在选定某个 $^L\neg\beta$ 之后，
我们希望把 $\MH({}^L\neg\alpha,{}^L\neg\beta,{}^L\neg\gamma;p)$ 上的一个膜
与给定的平坦丛 $E\to
C\backslash p$ 联系起来。
若 $V$（从而 $U$）是正则且半单的，则没有问题。
$\MH({}^L\neg\alpha,{}^L\neg\beta,{}^L\neg\gamma;p)$
是光滑的（对任意 ${}^L\neg\beta$），
且可取 ${\rm Im}\,\tau$
充分大而半经典地处理它。平坦丛 $E$ 对任意选定的 ${}^L\neg\beta$ 与 ${}^L\neg\eta$
都决定一个零膜 ${\mathcal B}_E$。

由于 $U$ 半单，第~\ref{sec:4.5}~节所研究的单值性
是阿贝尔的，其作用是与线丛
${\mathcal L}^v$ 对膜的 Chan-Paton 丛作张量积。
但限制到一个点上时 ${\mathcal L}^v$ 平凡。所以 ${\mathcal B}_E$
是阿贝尔单值性群的本征膜。它同样是作用于一般点 $q\in C$
的威尔逊线算符的本征膜。这可以按
无分歧时的同样推理得到；见文献~\cite{ref8} 第 8.2 节。作为
所有相关运算的本征膜，${\mathcal B}_E$
是无分歧时支撑在光滑点的零膜的良好类比。

总之，对正则半单的 $V$，我们处于与
无分歧时本质上相同的情形。
现在设 $V$ 不再是半单的，但仍然正则。
应记住的一个例子是：$V$ 可能是 $^LG_\C$ 的一个正则幂幺元，
例如 $SL(N,\C)$ 中的元素

% 式 (4.20)
\setcounter{equation}{19}
\begin{equation}\label{eq:4.20}
  V=\left(\begin{matrix}1& 1& 0&\dots&0\\
                         0&1&1&\dots&0\\
                           & & &\vdots&\\
                           0&0&0&\dots&1\end{matrix}\right)
\end{equation}

若 $V$ 不是半单的，则序对 $({}^L\neg\alpha,{}^L\neg\gamma)$
不正则；
例如，若 $V$ 幂幺，则 ${}^L\neg\alpha={}^L\neg\gamma=0$。
此时 $\MH({}^L\neg\alpha,{}^L\neg\beta,\gamma;p)$ 视 ${}^L\neg\beta$ 之不同，
可能带有第~\ref{sec:3.6}~节与第~\ref{sec:3.8}~节所述的局部奇点。
这恰好发生在三元组
$({}^L\neg\alpha,{}^L\neg\beta,{}^L\neg\gamma)$ 不正则之时。

\begin{yizhu}[原文疑误]上式中的 $\MH({}^L\neg\alpha,{}^L\neg\beta,\gamma;p)$ 缺了 $\gamma$ 上的左标 $^L$：
本段讨论的是 $^LG$ 一侧的参数，前一句刚写了 $({}^L\neg\alpha,{}^L\neg\gamma)$，紧接的三元组也写作
$({}^L\neg\alpha,{}^L\neg\beta,{}^L\neg\gamma)$，故此处应为 ${}^L\neg\gamma$。照原文保留，建议核对 PDF 第 107 页。\end{yizhu}

即便 $\MH({}^L\neg\alpha,{}^L\neg\beta,{}^L\neg\gamma;p)$
带有局部奇点，只要 $V$ 正则，点 $x_E$ 仍远离
该奇点。这意味着定义膜 ${\mathcal B}_E$ 没有问题。由于它以自然且唯一的方式
定义，它也是所有单值性运算的本征膜。所以对正则的 $V$，
我们再次处于与无分歧时本质上相同的情形。

\begin{yizhu}[易错点]本节的关键分层是“半单”“正则但非半单”“非正则”三档，切勿把前两者混为一谈：
$V$ 正则但非半单（如正则幂幺元）时，$({}^L\neg\alpha,{}^L\neg\gamma)$ 已不正则（$V$ 幂幺时二者均为零），
$\MH$ 可以有局部奇点，但 $x_E$ 不在奇点轨迹上，${\mathcal B}_E$ 仍有典范定义；只有当 $V$ 本身非正则时，
$x_E$ 才落入局部奇点轨迹，此时才需用 ${}^L\neg\beta$ 爆破出正维数的 $\Upsilon_E$，并只能给出膜的“范畴”。\end{yizhu}

若 $V$ 非正则，就会出现问题。此时若取
${}^L\neg\beta=0$，点 $x_E$ 便落在局部奇点的
轨迹之中。这给定义膜 ${\mathcal B}_E$ 带来困难。
我们可以取 ${}^L\neg\beta\not=0$ 且足够一般，
使三元组 $({}^L\neg\alpha,{}^L\neg\beta, {}^L\neg\gamma)$ 正则，从而消解奇点。
${}^L\neg\beta$ 是复结构 $J$ 中的 Kähler 参数，
取足够一般的 ${}^L\neg\beta$ 的效果是把局部奇点的轨迹
爆破掉。（这给出局部上幂零锥的施普林格消解
的整体版本。）这使 $\MH({}^L\neg\alpha,{}^L\neg\beta,{}^L\neg\gamma)$
非奇异，但爆破把点 $x_E$ 换成了
一个正维数的代数簇 $\Upsilon_E$。我们能做的最好的事，
就是把整个支撑在 $\Upsilon_E$ 上的膜的族或“范畴”
与 $E$ 联系起来。单值性群作用在这个
范畴上，其作用方式可能在某种意义上是不可约的。

另一种做法是保持 ${}^L\neg\beta=0$（或
取得足够特殊，以避免爆破点 $x_E$），而
通过取足够一般的 ${}^L\neg\eta$ 使西格玛模型光滑。这样选取
${}^L\neg\eta$ 或许能使我们定义一个
零膜 ${\mathcal B}_{E,{}^L\neg\eta}$。但足够一般的
${}^L\neg\eta$ 的可能选取被分成若干“房（chamber）”。我们猜测，若可以用某个
${}^L\neg\eta$ 的选取来定义零膜 ${\mathcal B}_{E,{}^L\neg\eta}$，则该
零膜依赖于包含 ${}^L\neg\eta$ 的
“房”。

例如，在幂幺单值性即 ${}^L\neg\alpha={}^L\neg\gamma=0$ 的情形，
若同时保持 ${}^L\neg\beta=0$ 以避免消解奇点，则
使西格玛模型不出现奇点对 ${}^L\neg\eta$ 的条件
是 ${}^L\neg\eta$ 必须正则。但正则 ${}^L\neg\eta$ 构成的空间
并不连通，因为非正则的 ${}^L\neg\eta$ 把 ${}^L\neg\TT$
分成若干仿射外尔房。
若我们试图通过变动 ${}^L\neg\beta$ 来从一个房
过渡到另一个房，以便绕过分隔不同房的奇点，
就会遇到用来定义群 ${\mathcal V}_U$ 的那些
单值性。

结果是：对非正则的 $V$，我们可以在希格斯丛模空间上定义一族
或一个“范畴”的膜，
它们与平坦丛 $E\to C\backslash p$ 相联系，并支撑在 $x_E$ 或其爆破之上。
这个范畴被非阿贝尔
群 ${\mathcal V}_U$ 作用。该范畴中所有膜都是作用于
远离 $p$ 之处的威尔逊算符的本征膜。但显然没有
办法在这族膜中挑出某个特定的膜。
也许 ${\mathcal V}_U$ 在某种意义上不可约地作用
在这族膜上。

应用对偶变换 $S:\tau\to -1/n_{\mathfrak g}\tau$，
在规范群为 $G$ 的 $A$-模型中得到同样的图景：一族 $A$-膜，
它们被 $\mathcal V_U$ 作用，并且是作用于 $p$ 之外各点的
't Hooft／赫克算符的本征膜。最后，利用 c.c. 膜，可以把这族膜映为
$\M(\alpha,p;G)$ 上的一族扭曲 ${\mathcal D}$-模，它们以同样的方式
与赫克算符及群 $\mathcal V_U$ 相联系。

因此，对非正则的 $V$，几何朗兰兹纲领可以表述为一个对偶：
它把一族被 $\mathcal V_U$ 作用的 $B$-膜，映为
一族被同一个群作用的
$A$-膜或 ${\mathcal D}$-模。

\begin{yizhu}[含糊点拨]此处“只能得到范畴、得不到典范零膜”的论证作者未展开，实质是：
非正则 $V$ 迫使 ${}^L\neg\beta$ 取一般值以爆破局部奇点，而爆破把点 $x_E$ 换成正维数的 $\Upsilon_E$；
$\Upsilon_E$ 上并无典范的零膜，且 $\mathcal V_U$ 的作用一般不可约，故只能留下整族膜。反过来，
正合列~\eqref{eq:4.15} 说明 $B_Y$ 一般严格大于 $\mathcal V_Y$，正是这一“额外的 $\pi_1$”使把线算符
唯一提升为单值性变换成为不可能。\end{yizhu}

\bigskip\noindent{\it 寻找典范零膜}

我们刚刚看到：用与极大环面 $\TT\subset G$ 相联系的
标准面算符所能推出的对偶陈述，
在带分歧的平坦丛 $E$ 具有非正则单值性 $V$ 时会变得更
复杂。尽管这一复杂性大概会让走到这一步的物理学家（若有的话）吃惊，
对几何学家却毫不意外，因为根据与数论的类比，这本就在预料之中。

不过，我们可以提出如下问题。给定带非正则单值性 $V$ 的平坦丛 $E\to C\backslash p$，
能否找到关于这个平坦丛的另一个对偶陈述，其中涉及的是
典范零膜而非整个范畴？若答案是“否”，
并不清楚这有什么不妥；但若答案是“是”，我们愿意把它找出来。

对 $V$ 的某些共轭类，我们可以简单地用基于一般列维子群 $\mathbb{L}$
的面算符（而非迄今所用的一般面算符）来回答这个问题。
在存在 $\mathbb{L}$ 型面算符时，希钦方程的一个解
在复结构 $J$ 中描述一个平坦丛 $E\to C\backslash p$，
其可能的单值性 $V$ 已在第~\ref{sec:3.3}~节末尾分析过。这个单值性
共轭于 $\EUP$ 中的一个元素，而 $\EUP$ 是由序对 $(\mathbb{L},{}^L\neg\alpha)$ 决定的
抛物子群。
一般地，该单值性是 $\mathbb{L}$ 正则的（其意义见第~\ref{sec:3.3}~节）。
恰在 $V$ 为 $\mathbb{L}$ 正则时，带分歧的平坦丛
$E$ 对应于 $\M_{H,\mathbb{L}}$ 上的一个光滑点，
从而对应于靶空间为 $\M_{H,\mathbb{L}}$ 的西格玛模型中的典范零膜
${\mathcal B}_{E}$。
按与上面相同的论证，这个典范零膜是所有相关威尔逊算符与
单值性的本征膜。对它应用 $S$-对偶，便得到辛结构为 $\omega_K$ 的
$A$-模型（规范群为 $G$ 而非 $^L\neg G$）中的一个膜，
然后利用 $A$-膜与 ${\mathcal D}$-模之间的关系，我们把 $E$ 与
抛物丛模空间 ${\mathcal M}_\mathbb{L}(\alpha;p,G)$ 上的一个 ${\mathcal D}$-模联系起来。

总之，若单值性对某个 $\mathbb{L}$ 是 $\mathbb{L}$ 正则的，
则通过考虑带 $\mathbb{L}$ 型面算符的理论，我们可以化归到与
无分歧情形一样简单的情形。
例如，回到式~\eqref{eq:3.30} 所描述的、$G=SL(3,\C)$ 的列维子群，
$\mathbb{L}$ 正则共轭类的典型例子是：包含元素
${\rm diag}(\lambda,\lambda,\lambda^{-2})$（$\lambda^3\not=1$）的半单共轭类，
以及代表元由式~\eqref{eq:3.32} 给出的 Richardson 共轭类。

对 $G=SL(N,\C)$，每个共轭类都对某个 $\mathbb{L}$ 是 $\mathbb{L}$ 正则的。
给定 $V\in SL(N,\C)$，
只需取 $U\in \TT$ 落在 $V$ 的轨道的闭包中，再令 $\mathbb{L}$
为 $SL(N,\C)$ 中与 $U$ 交换的子群。然而，对其他群，
情况并非如此。例如，$SL(N,\C)$ 之外的单李群含有
刚性的非中心轨道。（若一条轨道不能形变到任何邻近的
轨道，就称它是刚性的；通常是因为所有邻近轨道都有更大的维数。）
这样的轨道对任何 $\mathbb{L}$ 都不是 $\mathbb{L}$ 正则的。

我们猜测，有可能在 ${\mathcal N}=4$ 超杨–米尔斯理论中定义具有如下性质的
更多超对称面算符。
对每个带分歧的平坦丛 $E\to C\backslash p$，
我们希望能选到一个面算符，
使得在存在该面算符时，$E$ 对应于 $\Psi=\infty$ 处 $B$-模型中的
典范零膜 ${\mathcal B}_E$。此外，这个面算符应当
有 $S$-对偶（它将是一个大致同类型的面算符）。
把 $S$-对偶应用于零膜 ${\mathcal B}_E$，
将在 $\Psi=0$ 处的对偶模型中给出一个 $A$-膜。
然而，这些新的面算符不会与抛物
子群相关，所以在存在这种算符时，我们不
指望 $\MH$ 以通常方式与余切丛相联系。因此，
就不可能把对偶的 $A$-膜与 ${\mathcal D}$-模联系起来，
而对偶将不得不表述为 $\MH({}^L\neg G)$ 的 $B$-模型
与 $\MH(G)$ 的 $A$-模型之间的镜像对称或 $S$-对偶，
而不是 $\MH$ 上的相干层（或 $B$-膜）与某个别的空间上的
${\mathcal D}$-模之间的关系。
缺少通过 ${\mathcal D}$-模的解释，或许意味着这个更一般的
对偶（若存在）与数论无关，但它仍可能给出
$S$-对偶与镜像对称的一个漂亮应用。

\begin{yizhu}[物理图像补充]把 §4.5--§4.6 的图景串起来看：无分歧时，威尔逊／'t Hooft 线算符作用在一般点 $q$，
相当于在 $q$ 处插入一个表示（或 Hecke 修正）；当 $q$ 趋近歧点 $p$ 时，这一作用退化为分歧参数的单值性群
$\mathcal V_Y$（或 $B_Y$）对膜的“群半环”作用。于是“$p$ 处取代线算符的结构”就是歧点处的
Hecke 修正，而 $B$-膜的范畴（$J$ 型 $B$-模型）经 $S$-对偶与 c.c. 膜对应到 $A$-膜及
$\M$ 上的扭曲 ${\mathcal D}$-模范畴。非正则 $V$ 之所以使“典范零膜”消失，正是因为
Hecke 修正不再有唯一选择，只剩范畴——这正是数论中非半单 Langlands 参数没有唯一对应物的几何对应。\end{yizhu}
